%!TEX root = ../main.tex
% =============================================================
%  CAPÍTULO 8 — IMPLICANCIAS, CONCLUSIONES Y RECOMENDACIONES
%                                       (Entrega 8 · 5 % · 20/10)
% =============================================================

\chapter{Implicancias, conclusiones y recomendaciones}\label{cap:conclusiones}

\begin{guia}[Guía de la cátedra: qué se evalúa en este capítulo]
\begin{enumerate}
    \item \textbf{Respuesta explícita a la hipótesis} del
          capítulo~\ref{cap:introduccion}, con matices y contraejemplos
          si corresponde (``afirmativa, con matices'' es una respuesta
          válida; ``se cumplieron todos los objetivos'' sin evidencia no).
    \item \textbf{Implicancias}: técnicas, sociales y económicas.
    \item \textbf{Limitaciones honestas}, por categoría (técnicas, de datos,
          de validación, de mercado), sin maquillaje. Declarar como
          preliminar lo que es preliminar.
    \item \textbf{Recomendaciones} accionables y \textbf{trabajo futuro}.
    \item \textbf{Contribución al conocimiento}: qué sabe ahora la disciplina
          que antes no sabía.
\end{enumerate}
Cierra el arco narrativo: cada objetivo específico se retoma y se
indica en qué medida se alcanzó, con referencia a la sección que lo
demuestra.
\end{guia}

Este capítulo cierra el trabajo con la evidencia disponible al 15 de
octubre de 2026, la misma del capítulo~\ref{cap:resultados}. El piloto
con usuarios reales se desarrolla entre el 27 de octubre y el 27 de
noviembre de 2026, después de la entrega de este capítulo. Por ese
motivo, y tal como anticipó el riesgo R-03 del
capítulo~\ref{cap:riesgos}, la respuesta a la hipótesis es preliminar y
se completará con los indicadores del piloto en la entrega final.

\section{Conclusiones}\label{sec:conclusiones}

\subsection{Hallazgos principales}

\begin{enumerate}
    \item \textbf{El sistema existe y está desplegado.} Se construyó un
          \gls{mvp} de nueve módulos, con 160 requerimientos funcionales,
          37 registros de decisión de arquitectura y 739 pruebas
          automatizadas en el backend, con una cobertura de líneas del
          86,2\,\%, desplegado en \path{tinku.site}
          (capítulo~\ref{cap:resultados}). Lo construyó una sola persona,
          con un método dirigido por especificaciones y asistido por
          \gls{ia}.
    \item \textbf{Las pruebas automatizadas no bastaron para garantizar la
          seguridad.} La auditoría técnica encontró 36 hallazgos, siete
          de ellos críticos y vinculados a la protección del menor, en un
          sistema cuyas 383 pruebas de ese momento resultaban exitosas.
          Los 36 se cerraron con una prueba de regresión. Las pruebas
          sobre el despliegue real encontraron después tres hallazgos de
          severidad alta que la suite tampoco detectaba, uno de ellos un
          hueco de la propia especificación.
    \item \textbf{El problema y la modalidad tienen respaldo; el modelo de
          cobro, no del todo.} La encuesta ($n = 88$) respalda el
          diagnóstico de informalidad, la aceptación de la tutoría virtual
          (77,3\,\%) y el valor de la verificación de los tutores
          (81,8\,\%). En cambio, solo el 35,2\,\% prefiere pagar dentro de
          una plataforma, y la mediana de la disposición a pagar es la
          mitad del precio publicado en Superprof con el que se calcula el
          ticket (riesgos R-02 y R-05). La docente entrevistada coincide en
          el valor del apoyo particular y en la equivalencia de la clase
          virtual, pero difiere en el pago: lo preferiría dentro de la
          plataforma ``por seguridad''.
    \item \textbf{La \gls{ia} se valora por lo que resuelve, no por su
          nombre.} El 68,2\,\% quiere que la plataforma le recomiende un
          tutor, pero solo el 5,7\,\% eligió la ``\gls{ia} que recomiende
          profesores'' como función importante. El resumen automático,
          descrito por su utilidad, fue la segunda función más valorada
          (56,8\,\%). La entrevistada pidió que la recomendación se
          explique por la reputación del tutor y por los resultados de
          alumnos similares.
    \item \textbf{La viabilidad económica es condicional.} Con la
          comisión del \comision{}, el escenario base resulta viable
          (\gls{van} de \usd{18.869} y \gls{tir} del 41,57\,\%), igual que
          el optimista (\gls{van} de \usd{51.990}), pero el pesimista no
          lo es (\gls{van} de \usd{-4.410}). La viabilidad depende de
          sostener unas 5.858 sesiones anuales en promedio, de no bajar la
          comisión de aproximadamente 18,2\,\% y de operar sin un equipo
          rentado (capítulo~\ref{cap:economica}).
    \item \textbf{Contrastar el diseño con el sistema cambió las
          decisiones.} La contrastación del modelo económico con el
          sistema implementado mostró que la comisión de MercadoPago la
          paga el tutor y no la plataforma, y cambió la conclusión del
          capítulo~\ref{cap:economica}. La imposibilidad de anonimizar el
          audio antes de enviarlo a un proveedor externo llevó a
          transcribir primero, enviar al \gls{llm} solo el texto y excluir
          el resumen de las clases con menores (ADR-M6-03).
\end{enumerate}

\subsection{Respuesta a la hipótesis}\label{sec:respuesta-hipotesis}

La hipótesis central sostiene que una plataforma con \gls{matching}
semántico permitirá a estudiantes de zonas con escasez de tutores
acceder a apoyo académico de calidad equivalente a la de los centros
urbanos. \textbf{Con la evidencia disponible, la hipótesis no puede
confirmarse ni refutarse: la evidencia es consistente con sus
condiciones de posibilidad, pero todavía no verifica su núcleo.}

La evidencia respalda tres condiciones necesarias para que la hipótesis
se cumpla: la plataforma es técnicamente construible y operable por un
equipo mínimo; la modalidad virtual es aceptada por el 77,3\,\% de los
encuestados y considerada pedagógicamente equivalente por la docente
entrevistada; y el modelo económico es viable en los escenarios base y
optimista. En cambio, falta evidencia sobre el núcleo de la hipótesis: el
acceso efectivo de estudiantes del interior. El 90,9\,\% de las
respuestas de la encuesta proviene de la provincia de Buenos Aires y de
la Ciudad Autónoma de Buenos Aires, la entrevistada trabaja en Pilar y
no conoce casos concretos de la brecha rural, y el piloto todavía no
comenzó.

Las dos afirmaciones técnicas del capítulo~\ref{cap:introduccion} quedan
en el siguiente estado:

\begin{enumerate}
    \item \textbf{Relevancia de las recomendaciones del \gls{matching}.}
          Pendiente. El motor está implementado, y la calibración con 40
          consultas de elaboración propia muestra que separa las
          relevantes de las irrelevantes: con el umbral elegido, el 95\,\%
          de las relevantes conserva resultados y ninguna irrelevante los
          obtiene. La aceptación por parte de los usuarios (indicador
          I-01, umbral superior al 60\,\%) solo puede medirse en el piloto.
          \completar{Resultado de I-01 en el piloto.}
    \item \textbf{Sesiones entre provincias con calidad aceptable.}
          Pendiente. El aula virtual funcionó con salas de video reales en
          las pruebas en producción, pero la calidad percibida (I-02) y la
          sesión entre provincias con una del interior (I-03) se miden en
          el piloto.
          \completar{Resultados de I-02 e I-03 en el piloto.}
\end{enumerate}

La respuesta, por lo tanto, es \textbf{preliminar}: la hipótesis se
sostiene como plausible, pero su componente geográfico, que es el que
distingue a \sistema{} de las plataformas existentes, carece hoy de
evidencia directa.

\subsection{Cumplimiento de los objetivos específicos}

La tabla~\ref{tab:cumplimiento} retoma cada objetivo específico del
capítulo~\ref{cap:introduccion} y el grado en que se alcanzó.

\begin{table}[H]
    \centering
    \caption{Grado de cumplimiento de los objetivos específicos}
    \label{tab:cumplimiento}
    \small
    \begin{tabularx}{\textwidth}{>{\raggedright\arraybackslash}p{2.6cm} L l L}
        \toprule
        \tb{Objetivo} & \tb{Evidencia} & \tb{Grado} & \tb{Observación} \\
        \midrule
        OE1 Mercado y estado del arte & Capítulo~\ref{cap:marco-teorico}; apéndice~\ref{apx:encuesta} & Alcanzado & Encuesta concentrada en Buenos Aires; una sola entrevista realizada, la del Perfil~2 está pendiente (apéndice~\ref{apx:entrevista}) \\
        OE2 Arquitectura & Capítulo~\ref{cap:marco-tecnologico}; apéndice~\ref{apx:uml} & Alcanzado & --- \\
        OE3 Motor de \gls{matching} & Capítulo~\ref{cap:resultados}, pruebas en producción & Parcial & Implementado y calibrado con consultas propias; sin validar con usuarios (I-01) \\
        OE4 Aula virtual & Capítulo~\ref{cap:resultados}, pruebas en producción & Parcial & Aula y kill-switch funcionaron de punta a punta en la prueba en vivo; faltan la exactitud del clasificador con contenido real y la latencia en dispositivos reales, por lo que las clases con menores siguen deshabilitadas (R-01) \\
        OE5 Pagos & Capítulo~\ref{cap:resultados}; capítulo~\ref{cap:economica} & Parcial & Implementado con \gls{oauth} por tutor; pendiente de verificar en el entorno de pruebas de MercadoPago (R-11) \\
        OE6 \gls{ia} educativa & Capítulo~\ref{cap:resultados}, listado~\ref{lst:puerto-resumen} & Parcial & Implementado; apagado en producción hasta contar con la cláusula legal de grabación (R-16) \\
        OE7 Reputación & Capítulo~\ref{cap:resultados}, módulo de reputación & Alcanzado & Umbral mínimo de calificaciones antes de mostrar el promedio \\
        OE8 Validación funcional & Capítulo~\ref{cap:resultados}, pruebas de aceptación & Pendiente & Piloto del 27 de octubre al 27 de noviembre de 2026, sin clases con menores \\
        \bottomrule
    \end{tabularx}
\end{table}

\subsection{Interpretación frente a los estudios previos}

Los resultados dialogan con la literatura del
capítulo~\ref{cap:marco-teorico} en cuatro puntos.

\begin{itemize}
    \item \textbf{Efectividad de la tutoría.} La encuesta confirma que el
          apoyo particular se busca y se valora, y la entrevistada lo
          considera ``fundamental'' porque la escuela no puede seguir a
          cada alumno. Ambos resultados son coherentes con
          el problema de los 2 sigma \citep{bloom1984} y con el
          meta-análisis de tutoría privada \citep{sciencedirect2022}. Este
          trabajo no aporta evidencia nueva sobre el efecto de la tutoría
          en el aprendizaje, ya que no midió rendimiento académico.
    \item \textbf{\Gls{matching} semántico.} Los antecedentes
          \citep{isotropy2026, paperrec2025} validaron la técnica con
          pruebas de relevancia percibida. La encuesta agrega un matiz que
          esos trabajos no abordan: el usuario valora la recomendación,
          pero no la etiqueta ``\gls{ia}''. La relevancia técnica de la
          recomendación sería, así, necesaria pero no suficiente; también
          importa que pueda explicarse.
    \item \textbf{Reputación bidireccional.} El diseño adoptó la
          reputación bidireccional discutida por \citet{reputation}. La
          entrevista introduce una variable que esa discusión no
          contempla: cuando una de las partes es menor de edad, quién
          califica deja de ser un detalle de interfaz y pasa a ser una
          decisión de protección.
    \item \textbf{Mercados bilaterales.} El problema de arranque en frío
          \citep{rochet2006} aparece en los datos: solo 9 de las 88
          respuestas corresponden a tutores, y la oferta es el recurso más
          escaso para iniciar el piloto (R-07). La comisión agrava el
          problema: con el \comision{}, el tutor recibe el 66,96\,\% de su
          precio, seis puntos menos que cuando la plataforma absorbía la
          pasarela.
\end{itemize}

\subsection{Limitaciones}\label{sec:limitaciones-finales}

\begin{itemize}
    \item \textbf{De validación.} La hipótesis se responde sin los datos
          del piloto, que se realiza después de esta entrega. El piloto,
          además, no incluye clases con menores, aunque los menores son una
          parte central de la población objetivo.
    \item \textbf{De datos.} La encuesta no es probabilística y el
          90,9\,\% de sus respuestas proviene de Buenos Aires. La
          entrevista es una sola, a una docente de Educación Física sin
          experiencia en clases particulares, y en algunas preguntas el
          resultado de la encuesta se presentó antes de la respuesta; la
          entrevista a una tutora particular, que aporta el lado de la
          oferta, está pendiente.
          La calibración del umbral de relevancia usa consultas de
          elaboración propia.
    \item \textbf{Técnicas.} El clasificador del kill-switch no se midió
          con contenido real ni en dispositivos reales; el resumen está
          apagado en producción; el cobro con \gls{oauth} no se verificó
          en el entorno de pruebas de MercadoPago; las pruebas por rol en
          producción deben repetirse tras la remediación (R-15) y el
          despliegue debe volver a su estado previo antes del piloto
          (R-17). Todo el sistema corre en un único servidor propio
          (R-14).
    \item \textbf{De mercado.} La disposición a pagar se relevó sobre una
          muestra urbana y resultó la mitad del precio con el que se
          calcula el ticket; el precio y la demanda del interior pueden
          diferir. La preferencia por pagar dentro de la plataforma es
          minoritaria.
\end{itemize}

\section{Implicancias}\label{sec:implicancias}

\begin{itemize}
    \item \textbf{Teórica.} El trabajo sugiere que, en un
          \gls{marketplace} educativo, la aceptación del \gls{matching} no
          depende solo de la relevancia técnica de la recomendación, sino
          de que el usuario la perciba por su resultado. También muestra
          que la presencia de menores modifica decisiones que la
          literatura trata como neutras, como qué datos pueden enviarse a
          un proveedor de \gls{ia} o quién paga la clase.
    \item \textbf{Práctica.} La combinación de especificaciones previas,
          registros de decisión y asistencia de \gls{ia} permitió que una
          sola persona construyera y desplegara un sistema de nueve
          módulos. La auditoría y las pruebas en producción mostraron el
          límite de ese método: la velocidad de construcción no reemplaza
          la revisión independiente, porque las pruebas verifican lo que
          se especificó y no lo que se omitió.
    \item \textbf{De innovación.} Ninguna de las plataformas relevadas
          integra \gls{matching} semántico, aula virtual, pago con
          \gls{escrow} y controles específicos para menores en un mismo
          flujo (tabla~\ref{tab:comparativa}). La integración es la
          innovación; ninguna pieza por separado es nueva.
    \item \textbf{Social y económica.} Si el piloto confirma la hipótesis,
          la plataforma ampliaría el acceso al apoyo escolar en zonas sin
          oferta presencial y formalizaría ingresos de tutores hoy
          informales. La condición es económica: por debajo de unas 1.774 a
          2.194 sesiones anuales la operación no cubre sus costos fijos, y
          un servicio que cierra no reduce ninguna brecha.
\end{itemize}

\section{Aportes}\label{sec:aportes}

\begin{enumerate}
    \item Un \gls{mvp} funcional, desplegado y documentado, con
          especificaciones, 37 registros de decisión de arquitectura y 739
          pruebas en el backend, que otra persona puede retomar
          (anexo~\ref{anx:repositorios}).
    \item Un modelo de protección del menor para tutoría virtual 1 a 1:
          cuenta pagadora del adulto responsable, \gls{cap} obligatorio
          para los tutores de menores, kill-switch que corre en el
          navegador de cada participante y exclusión del resumen
          automático en las clases con menores.
    \item Un modelo económico reproducible, contrastado con el sistema
          implementado, que muestra cómo esa contrastación cambió la
          conclusión (anexo~\ref{anx:modelo-financiero}).
    \item Evidencia empírica de que, en este mercado, la \gls{ia} se valora
          por su resultado y no como atributo del producto.
\end{enumerate}

\section{Recomendaciones}\label{sec:recomendaciones}

\begin{itemize}
    \item \textbf{Para la práctica profesional.} Complementar las pruebas
          automatizadas con una revisión independiente y con pruebas sobre
          el despliegue real antes de cada pasaje a producción, en
          especial en los controles de seguridad. En el desarrollo
          asistido por \gls{ia}, verificar lo que el sistema hace y no solo
          lo que la documentación declara.
    \item \textbf{Para los investigadores.} Medir la aceptación de un
          sistema de recomendación con y sin explicación, y separar en las
          encuestas la valoración de una función de la valoración de la
          tecnología que la resuelve.
    \item \textbf{Para instituciones y políticas públicas.} Ofrecer una
          vía de verificación de antecedentes consultable por las
          plataformas, con consentimiento del titular, que reemplace la
          revisión manual de un certificado presentado por el propio
          tutor (R-10). Las escuelas podrían, además, derivar alumnos y
          aportar el volumen que el modelo económico necesita (O-02).
    \item \textbf{Para el desarrollo tecnológico de \sistema{}.} Antes del
          piloto, repetir las pruebas por rol en producción y revertir el
          estado del despliegue (R-15, R-17). Antes de habilitar las
          clases con menores, medir la exactitud del clasificador y su
          latencia en dispositivos reales (R-01). Antes de operar
          comercialmente, verificar el cobro en el entorno de pruebas de
          MercadoPago (R-11), revisar la comisión con los datos del piloto
          y mostrar el motivo de cada recomendación. Una vez realizado el
          piloto, reconsiderar las propuestas de la entrevista que no se
          incorporaron al \gls{mvp}: la presencia del adulto responsable en
          las clases con menores y la calificación del tutor por parte del
          adulto.
\end{itemize}

\section{Futuras líneas de investigación}\label{sec:futuras-lineas}

\begin{itemize}
    \item \textbf{Validación en el interior.} Replicar la encuesta y el
          piloto en localidades con escasa oferta de tutores, que es la
          población a la que apunta la hipótesis (R-08).
    \item \textbf{Resumen automático con menores.} Evaluar la
          transcripción local o en infraestructura propia, previa a la
          anonimización, que permita ofrecer el resumen sin transferir la
          voz del menor a un proveedor externo.
    \item \textbf{Efecto en el aprendizaje.} Medir el rendimiento académico
          de los alumnos antes y después de un ciclo de tutorías, para
          contrastar el resultado con la literatura sobre tutoría.
    \item \textbf{Continuidad entre clases.} Estudiar el valor de las
          tareas y entregas asincrónicas propuestas en la entrevista.
    \item \textbf{Funciones excluidas del alcance.} Aplicaciones móviles
          nativas, integración con sistemas estatales, accesibilidad
          completa, suscripción, clases grupales y el asistente
          conversacional de repaso (sección~\ref{sec:alcance}).
\end{itemize}
